\section{Cornerstone Analysis}
\label{sec:analysis}

To evaluate the Cornerstone octree construction method under realistic conditions, we developed a configurable batch generation pipeline (\texttt{run\_batches}) that sweeps over a matrix of particle distributions, particle counts, and perturbation magnitudes. Particles that leave the domain after perturbation are truncated back to the boundary. 

Much of the already existing performance analysis in regards to Cornerstone has been for extremely large particle counts with computation distributed across thousands of GPUs. For smaller application spaces such as real-time rendering or small simulations, it is possible that a low number of GPUs are either accessible or optimal for a given workload. Due to technical constraints, our preliminary results and analysis operate using only a single GPU with relatively low particle sizes. However, low particle count workloads involving only one GPU are not uncommon, and analysis on these workloads may yield useful conclusions for a number of applications.

\subsection{Distribution Generators}
All initial particle distributions are confined to the cube domain $[-1,1]^3$, where all distributions are as large as the domain.
Six initial particle distributions are used:

\begin{enumerate}
    \item \textbf{\texttt{uniform\_initial}}: Samples particles uniformly at random from the domain.
    \item \textbf{\texttt{pancake\_initial}}: A cylindrical distribution along the Z-axis where the Z-axis is strongly compressed, as to resemble a disk, frisbee, or (as the name suggests) pancake.
    \item \textbf{\texttt{pancake\_initial\_tilted}}: Same as the pancake but tilted off-axis, so the thin dimension is no longer aligned with a coordinate plane.
    \item \textbf{\texttt{filament\_z\_initial}}: A one-dimensional filament aligned along the $z$-axis, characterized by the parametric equation: (0, 0, t).
    \item \textbf{\texttt{filament\_yz\_initial}}: A filament lying in the $yz$-plane, anisotropic in two directions, characterized by the parametric equation: (0, t, t)..
    \item \textbf{\texttt{filament\_xyz\_initial}}: A filament oriented along the domain diagonal, anisotropic in all three coordinate directions simultaneously, characterized by the parametric equation: (t, t, t).
\end{enumerate}

By analyzing 1D, 2D, and 3D manifolds that are or are not axis-aligned, a wide variety of behavior can be observed after applying perturbations. Because the perturbation distribution is 3-dimensional, the various distributions will have their unique attributes degraded as the perturbation magnitude increases.  

\subsection{Perturbation Model}

After the initial distribution is generated, a perturbation vector is applied to simulate particle motion between timesteps. The perturbation function draws each displacement component independently from $\mathcal{U}(-s, +s)$ where $s$ is the scale factor. Particles that move outside the domain after perturbation are truncated to the boundary (\texttt{out\_of\_bounds = `truncate'}). Four perturbation scales are used:

\begin{itemize}
    \item $s = 0.001$: Very small perturbation; particles barely move relative to the domain.
    \item $s = 0.01$\phantom{0}: Moderate perturbation; noticeable displacement but locally coherent.
    \item $s = 0.1$\phantom{00}: Large perturbation; significant displacement triggering substantial tree restructuring.
    \item $s = 0.5$\phantom{00}: Very large perturbation; particles may traverse a large fraction of the domain.
\end{itemize}

\subsection{Runtime Configuration}

The Cornerstone framework accepts several runtime parameters via the command line:

\begin{itemize}
    \item \texttt{--bucket-size-global}: Maximum number of particles per leaf node in the global octree (default: 1024).
    \item \texttt{--bucket-size-focus}: Maximum number of particles per leaf node in the focused octree (default: 64).
    \item \texttt{--theta}: The multipole acceptance criterion (MAC) parameter controlling the accuracy/performance tradeoff (default: 0.6).
    \item \texttt{--gpu}: Flag to enable GPU-accelerated octree construction.
\end{itemize}

No rotations were applied (\texttt{rotations = None}).

\subsection{Test Configurations}

The full sweep comprises $6$ distributions $\times$ $5$ particle counts $\times$ $4$ perturbation scales $= 120$ configurations.

\subsection{Code}
We have created the necessary Jupyter notebooks and Python scripts to generate distributions of interest and their perturbations. This includes code to save those distributions to HDF5 files. We developed a framework in C++ for loading those HDF5 distributions and their perturbations and generating the corresponding octrees, adding the perturbations to the particles and then rebalancing the octree. This allows us to observe the effects of the initial construction and across varying magnitudes of perturbations. This includes methods for a CPU and GPU implementation as well as their related distributed MPI variants. This is flexible enough to allow us to generate any metrics of interest in the future and has already allowed us to gather the results discussed in \ref{subsec:results}. Tooling to instrument sections of interest is done using NVIDIA Tooling Extensions (NVTX). We instrumented our code to highlight important metrics at a relatively coarse level so that we could observe the primary steps necessary to building the Octree and focus our attention on bottlenecks. 

\subsection{Machines}
For the results used here we have used \href{https://www.cloudlab.us/}{CloudLab} to gain access to GPUs allowing us to get on a node with 4 V100s; however, access to this is difficult to attain on a regular basis and does not provide access to bare metal. We have access to the NCSA Delta cluster at this point and will use it to gather metrics on other GPUs and in varying configurations during the second half of the project. This should also provide much more consistent access to the resources we need. The Delta cluster comes with 4x and 8x GPU configurations per node with A40, A100, and H200 GPUs being the devices we are interested in testing on in the future. We expect that we will not do inter-node experiments as it is outside of the hardware limits we have.

% --- Uniform ---
% \begin{table*}[ht]
%     \centering
%     \caption{\texttt{uniform\_initial} results.}
%     \label{tab:uniform_results}
%     \begin{tabular}{|c|c|c|c|c|}
%         \hline
%         \textbf{$N$} & \textbf{Scale} & \textbf{Build (ms)} & \textbf{Rebalance (ms)} & \textbf{Depth} \\
%         \hline
%         $10^4$ & $0.001$ & & & \\
%         $10^4$ & $0.01$  & & & \\
%         $10^4$ & $0.1$   & & & \\
%         $10^4$ & $0.5$   & & & \\
%         \hline
%         $10^5$ & $0.001$ & & & \\
%         $10^5$ & $0.01$  & & & \\
%         $10^5$ & $0.1$   & & & \\
%         $10^5$ & $0.5$   & & & \\
%         \hline
%         $10^6$ & $0.001$ & & & \\
%         $10^6$ & $0.01$  & & & \\
%         $10^6$ & $0.1$   & & & \\
%         $10^6$ & $0.5$   & & & \\
%         \hline
%         $10^7$ & $0.001$ & & & \\
%         $10^7$ & $0.01$  & & & \\
%         $10^7$ & $0.1$   & & & \\
%         $10^7$ & $0.5$   & & & \\
%         \hline
%         $10^8$ & $0.001$ & & & \\
%         $10^8$ & $0.01$  & & & \\
%         $10^8$ & $0.1$   & & & \\
%         $10^8$ & $0.5$   & & & \\
%         \hline
%     \end{tabular}
% \end{table*}

% % --- Pancake ---
% \begin{table*}[ht]
%     \centering
%     \caption{\texttt{pancake\_initial} results.}
%     \label{tab:pancake_results}
%     \begin{tabular}{|c|c|c|c|c|}
%         \hline
%         \textbf{$N$} & \textbf{Scale} & \textbf{Build (ms)} & \textbf{Rebalance (ms)} & \textbf{Depth} \\
%         \hline
%         $10^4$ & $0.001$ & & & \\
%         $10^4$ & $0.01$  & & & \\
%         $10^4$ & $0.1$   & & & \\
%         $10^4$ & $0.5$   & & & \\
%         \hline
%         $10^5$ & $0.001$ & & & \\
%         $10^5$ & $0.01$  & & & \\
%         $10^5$ & $0.1$   & & & \\
%         $10^5$ & $0.5$   & & & \\
%         \hline
%         $10^6$ & $0.001$ & & & \\
%         $10^6$ & $0.01$  & & & \\
%         $10^6$ & $0.1$   & & & \\
%         $10^6$ & $0.5$   & & & \\
%         \hline
%         $10^7$ & $0.001$ & & & \\
%         $10^7$ & $0.01$  & & & \\
%         $10^7$ & $0.1$   & & & \\
%         $10^7$ & $0.5$   & & & \\
%         \hline
%         $10^8$ & $0.001$ & & & \\
%         $10^8$ & $0.01$  & & & \\
%         $10^8$ & $0.1$   & & & \\
%         $10^8$ & $0.5$   & & & \\
%         \hline
%     \end{tabular}
% \end{table*}

% % --- Pancake Tilted ---
% \begin{table*}[ht]
%     \centering
%     \caption{\texttt{pancake\_initial\_tilted} results.}
%     \label{tab:pancake_tilted_results}
%     \begin{tabular}{|c|c|c|c|c|}
%         \hline
%         \textbf{$N$} & \textbf{Scale} & \textbf{Build (ms)} & \textbf{Rebalance (ms)} & \textbf{Depth} \\
%         \hline
%         $10^4$ & $0.001$ & & & \\
%         $10^4$ & $0.01$  & & & \\
%         $10^4$ & $0.1$   & & & \\
%         $10^4$ & $0.5$   & & & \\
%         \hline
%         $10^5$ & $0.001$ & & & \\
%         $10^5$ & $0.01$  & & & \\
%         $10^5$ & $0.1$   & & & \\
%         $10^5$ & $0.5$   & & & \\
%         \hline
%         $10^6$ & $0.001$ & & & \\
%         $10^6$ & $0.01$  & & & \\
%         $10^6$ & $0.1$   & & & \\
%         $10^6$ & $0.5$   & & & \\
%         \hline
%         $10^7$ & $0.001$ & & & \\
%         $10^7$ & $0.01$  & & & \\
%         $10^7$ & $0.1$   & & & \\
%         $10^7$ & $0.5$   & & & \\
%         \hline
%         $10^8$ & $0.001$ & & & \\
%         $10^8$ & $0.01$  & & & \\
%         $10^8$ & $0.1$   & & & \\
%         $10^8$ & $0.5$   & & & \\
%         \hline
%     \end{tabular}
% \end{table*}

% % --- Filament Z ---
% \begin{table*}[ht]
%     \centering
%     \caption{\texttt{filament\_z\_initial} results.}
%     \label{tab:filament_z_results}
%     \begin{tabular}{|c|c|c|c|c|}
%         \hline
%         \textbf{$N$} & \textbf{Scale} & \textbf{Build (ms)} & \textbf{Rebalance (ms)} & \textbf{Depth} \\
%         \hline
%         $10^4$ & $0.001$ & & & \\
%         $10^4$ & $0.01$  & & & \\
%         $10^4$ & $0.1$   & & & \\
%         $10^4$ & $0.5$   & & & \\
%         \hline
%         $10^5$ & $0.001$ & & & \\
%         $10^5$ & $0.01$  & & & \\
%         $10^5$ & $0.1$   & & & \\
%         $10^5$ & $0.5$   & & & \\
%         \hline
%         $10^6$ & $0.001$ & & & \\
%         $10^6$ & $0.01$  & & & \\
%         $10^6$ & $0.1$   & & & \\
%         $10^6$ & $0.5$   & & & \\
%         \hline
%         $10^7$ & $0.001$ & & & \\
%         $10^7$ & $0.01$  & & & \\
%         $10^7$ & $0.1$   & & & \\
%         $10^7$ & $0.5$   & & & \\
%         \hline
%         $10^8$ & $0.001$ & & & \\
%         $10^8$ & $0.01$  & & & \\
%         $10^8$ & $0.1$   & & & \\
%         $10^8$ & $0.5$   & & & \\
%         \hline
%     \end{tabular}
% \end{table*}

% % --- Filament YZ ---
% \begin{table*}[ht]
%     \centering
%     \caption{\texttt{filament\_yz\_initial} results.}
%     \label{tab:filament_yz_results}
%     \begin{tabular}{|c|c|c|c|c|}
%         \hline
%         \textbf{$N$} & \textbf{Scale} & \textbf{Build (ms)} & \textbf{Rebalance (ms)} & \textbf{Depth} \\
%         \hline
%         $10^4$ & $0.001$ & & & \\
%         $10^4$ & $0.01$  & & & \\
%         $10^4$ & $0.1$   & & & \\
%         $10^4$ & $0.5$   & & & \\
%         \hline
%         $10^5$ & $0.001$ & & & \\
%         $10^5$ & $0.01$  & & & \\
%         $10^5$ & $0.1$   & & & \\
%         $10^5$ & $0.5$   & & & \\
%         \hline
%         $10^6$ & $0.001$ & & & \\
%         $10^6$ & $0.01$  & & & \\
%         $10^6$ & $0.1$   & & & \\
%         $10^6$ & $0.5$   & & & \\
%         \hline
%         $10^7$ & $0.001$ & & & \\
%         $10^7$ & $0.01$  & & & \\
%         $10^7$ & $0.1$   & & & \\
%         $10^7$ & $0.5$   & & & \\
%         \hline
%         $10^8$ & $0.001$ & & & \\
%         $10^8$ & $0.01$  & & & \\
%         $10^8$ & $0.1$   & & & \\
%         $10^8$ & $0.5$   & & & \\
%         \hline
%     \end{tabular}
% \end{table*}

% % --- Filament XYZ ---
% \begin{table*}[ht]
%     \centering
%     \caption{\texttt{filament\_xyz\_initial} results.}
%     \label{tab:filament_xyz_results}
%     \begin{tabular}{|c|c|c|c|c|}
%         \hline
%         \textbf{$N$} & \textbf{Scale} & \textbf{Build (ms)} & \textbf{Rebalance (ms)} & \textbf{Depth} \\
%         \hline
%         $10^4$ & $0.001$ & & & \\
%         $10^4$ & $0.01$  & & & \\
%         $10^4$ & $0.1$   & & & \\
%         $10^4$ & $0.5$   & & & \\
%         \hline
%         $10^5$ & $0.001$ & & & \\
%         $10^5$ & $0.01$  & & & \\
%         $10^5$ & $0.1$   & & & \\
%         $10^5$ & $0.5$   & & & \\
%         \hline
%         $10^6$ & $0.001$ & & & \\
%         $10^6$ & $0.01$  & & & \\
%         $10^6$ & $0.1$   & & & \\
%         $10^6$ & $0.5$   & & & \\
%         \hline
%         $10^7$ & $0.001$ & & & \\
%         $10^7$ & $0.01$  & & & \\
%         $10^7$ & $0.1$   & & & \\
%         $10^7$ & $0.5$   & & & \\
%         \hline
%         $10^8$ & $0.001$ & & & \\
%         $10^8$ & $0.01$  & & & \\
%         $10^8$ & $0.1$   & & & \\
%         $10^8$ & $0.5$   & & & \\
%         \hline
%     \end{tabular}
% \end{table*}
